% Metodologia -- Módulo 4 -- PODs
\begin{frame}{PODs: leitura operacional dentro dos regimes}
  \AipePipelineContext{4}
  \scriptsize\textbf{\textcolor{darkblue}{Syntetos-Boylan diz o regime macro; o POD traduz esse regime em um perfil operacional para decisão.}}
  \vspace{0.25em}

  \centering
  \renewcommand{\arraystretch}{1.08}
  \setlength{\tabcolsep}{3pt}
  \scriptsize
  \begin{tabular}{@{}p{0.28\linewidth}p{0.25\linewidth}p{0.35\linewidth}@{}}
    \toprule
    \textbf{Sinal observado} & \textbf{POD} & \textbf{Leitura operacional} \\
    \midrule
    muitos zeros e poucos picos & Sparse High Impact & risco concentrado em eventos raros \\
    ciclos ou picos recorrentes & High Vol. Seasonal & política sensível a períodos de alta \\
    mudança de nível & Unstable Trend & histórico recente ganha peso \\
    baixo volume estável & Low Vol. Stable & controle simples tende a bastar \\
    demanda frequente & Fast Moving & reposição contínua vira candidata \\
    \bottomrule
  \end{tabular}
  \vspace{0.08em}

  {\tiny\textcolor{slategray}{Fonte: capítulo de Metodologia, Tabela 4.1; demand\_profiling}}\\[-0.08em]
  {\tiny\textcolor{slategray}{Exemplo: no recorte BA, as séries são Lumpy no regime macro, mas os PODs separam subperfis úteis para a recomendação.}}
  \note{
    A mensagem do slide é separar duas ideias. Syntetos-Boylan classifica o
    regime macro da demanda a partir de ADI e CV ao quadrado. Os PODs não
    substituem essa taxonomia; eles refinam a leitura operacional usando
    sinais adicionais da série. Em vez de mostrar todos os limiares, a fala
    deve enfatizar a intuição: silêncios longos indicam Sparse High Impact,
    ciclos ou picos indicam High Vol. Seasonal, instabilidade indica
    Unstable Trend, baixo volume estável indica Low Vol. Stable, e demanda
    frequente indica Fast Moving. No Experimento 2, todas as séries são
    Lumpy no regime macro, mas os PODs ainda separam subperfis úteis. O POD
    organiza o contexto; a recomendação de política vem depois, pela
    simulação, geração de rótulos e PSE.
  }
\end{frame}
